\section*{第 \ref{ch:b3:scattering-theory} 章——散射理论}

\begin{solution}{exo:b3:scattering-theory:1}
(a) $nt = \qty{5.9e22}{m^{-2}}$。(b) $nt\sigma = \num{5.9e22}
\times \num{e-27} = \num{5.9e-5}$。(c) 散射率
$\num{6.2e12} \times \num{5.9e-5} = \qty{3.7e8}{s^{-1}}$；进入球面的
$\num{e-3}/4\pi$ 份额为 $\qty{2.9e4}{s^{-1}}$。(d) 所谓
“面积”同时体现相互作用与波长：同一个金
原子核对 α 粒子、电子和
中子呈现的“面积”各不相同，且随能量而变。
\end{solution}

\begin{solution}{exo:b3:scattering-theory:2}
(a) 两球中心距离小于 $R_1 + R_2$ 时发生碰撞：$\sigma = \pi(R_1 +
R_2)^2$。(b) $\sigma \approx \qty{2.8e-19}{m^2}$：$\ell = 1/n\sigma
\approx \qty{0.14}{\micro m}$——分子运动论的尺度（
$\sqrt2$ 因子由相对运动修正）。(c) 硬球
对每个瞄准距离环上的入射粒子，按角度均匀反射：
经典极限下各向同性。(d) 低能时的四倍
\omterm{def:b3:scattering-theory:cross-section}{截面}——纯粹的波动衍射，
任何经典阴影逻辑都不再适用。
\end{solution}

\begin{solution}{exo:b3:scattering-theory:3}
(a) 顶角为 $\theta$ 的等腰三角形：$q =
2k\sin(\theta/2)$。(b) $q$ 从 $0$ 变化到 $2k = \qty{2e15}{
m^{-1}}$：直至 $\sim\qty{5e-16}{m}$ 尺度的结构原则上均可编码
（实际上库仑势垒使 α 粒子
挡在外面）。(c) $k = E/\hbar c = \qty{2.5e15}{m^{-1}}$：分辨率
$\sim\qty{2e-16}{m}$——已深入质子内部。(d) 分辨率
$\sim 1/q_{\max} \sim \hbar/p$：提高能量，就是换来更短的
“标尺”。
\end{solution}

\begin{solution}{exo:b3:scattering-theory:4}
(a) $f$ 是长度；$|f|^2$ 是面积——按构造是
单位立体角上的量。(b) 将 $(\hbar k/m)|f|^2/r^2$ 在
球面上积分，得 $\Phi\sigma$。(c) 概率守恒：出射
波必须与前向束\emph{相消}干涉，以补偿
散射掉的部分；这一账目记在 $f$ 的相位中。(d)
光学定理正是这一“阴影核算”：总的
消散量 $=$ 前向干涉的亏损。
\end{solution}

\begin{solution}{exo:b3:scattering-theory:5}
(a) 角向积分给出 $4\pi\sin(qr)/qr$；再利用
$\int_0^\infty\eu^{-\mu r}\sin(qr)\,\dd r = q/(q^2 + \mu^2)$，
可得 $f = -2mg/\hbar^2(q^2 + \mu^2)$。(b) $\mu \to 0$：$|f|^2 =
(2mg/\hbar^2q^2)^2$；由 $q = 2k\sin(\theta/2)$ 与 $E =
\hbar^2k^2/2m$，这即卢瑟福公式。(c) 当 $\theta \to
0$ 时的发散积分为无穷：每个经过的粒子都会被
无限程力稍稍偏转；在物质中，原子
电子在超过 $\sim\qty{e-10}{m}$ 的距离上屏蔽原子核，并截断
这一发散。(d) $\hbar/m_\pi c = 197.3/140 = \qty{1.4}{fm}$
——核力的作用范围，由 π 介子质量预言。
\end{solution}

\begin{solution}{exo:b3:scattering-theory:6}
(a) 由 $\sin(ka + \delta_0) = 0$ 取最小解：$\delta_0
= -ka$。(b) $\sigma = (4\pi/k^2)\sin^2(ka) \to 4\pi a^2$。(c) 当
$\lambda \gg a$ 时，波把球视为点缺陷，
借助衍射在其四周重新合拢——“阴影”是一个短波概念，
而长波的结果计入整个球
表面 $4\pi a^2$。(d) 即使在短波下，
移除一圆盘的波也需要衍射填补——阴影
本身也会散射（阻挡的 $\pi a^2$ $+$ $\pi a^2$ 的
阴影衍射）。
\end{solution}

\begin{solution}{exo:b3:scattering-theory:7}
(a) $\pi(3 \times 5.4^2 + 23.7^2)\,\unit{fm^2} = \pi \times 649
\,\unit{fm^2} = \qty{20.4}{b}$——正是实测值。(b) 单态：
其巨大的 $|a_{\text{s}}|$ 压倒了其 $1/4$ 的权重。(c)
$a_{\text{t}}$ 大且为正：真实的束缚态（
氘核）刚好低于阈值；$a_{\text{s}}$ 大且
\emph{为负}：一个“虚”能级刚好高于阈值——几乎
几乎是第二个氘核，只是自然界没有将其束缚。(d) 氢的
原子核是中子最佳的动量匹配伙伴
（质量相等），每个质子二十靶恩的弹性散射，
使水成为极佳的紧凑型慢化剂。
\end{solution}

\begin{solution}{exo:b3:scattering-theory:8}
(a) $\sigma_0 = (4\pi/k^2)\sin^2\pi = 0$：$s$ 波恰好
如同不存在原子一样射出。(b) 波在
\emph{内部}被强烈扭曲，但出射时多出整数个
半波长：无可观测的相位偏移，无散射——
散射小波的相消干涉。(c) $KR
\approx \pi$，且 $E \ll V_0$：$V_0 \approx \pi^2\hbar^2/2mR^2
\approx \qty{9}{eV}$——原子尺度的势阱。(d) 它们
较浅的势阱在 eV 量级的
能量下从不使相位绕过 $\pi$。该效应经典上无法模拟：强吸引中的
透明窗口只对波存在。
\end{solution}

\begin{solution}{exo:b3:scattering-theory:9}
(a) $\sigma = 4\pi\sin^2\delta_0/k^2$。(b) $\operatorname{Im}f(0)
= \sin^2\delta_0/k$：$(4\pi/k)\operatorname{Im}f(0) = \sigma$——
已验证。(c) $\sin^2\delta_0 \le 1$：上限为 $4\pi/k^2 \approx
\qty{1e-20}{m^2} = \num{e8}$ 靶恩，对热中子而言——
余地仍高于最凶猛的吸收体（氙-135 的
百万靶恩）。(d) 上限即幺正性：波不能
移除多于干涉所允许的通量；强度使正弦
饱和，而永不超过它。
\end{solution}

\begin{solution}{exo:b3:scattering-theory:10}
(a) $\tau = \hbar/\Gamma = \num{6.58e-16}/0.027 =
\qty{2.4e-14}{s}$——按核的时间尺度很长：复合核
会“遗忘”其形成过程。(b) $4\pi/k^2 \approx
\num{3.9e5}$ b；观测值 $\num{2.3e4}$ b 为上限
乘以入射道的分支比（$\Gamma_n/\Gamma
\sim 0.06$）——共振以分宽度“售票”。(c)
$\Gamma/k_{\text{B}} \approx \qty{310}{K}$：加热燃料使
共振发生多普勒展宽，在慢化过程中捕获更多
中子——吸收随温度升高而增加，这是一种即时的
\emph{负}反馈，内建于铀本身。(d) 在
热能与快能之间，$^{238}$U 的共振“森林”吞噬
中子；燃料几何构型与慢化剂被设计成使中子
悄悄绕过它，其多普勒反馈则是反应堆
安全的支柱之一。
\end{solution}

\begin{solution}{exo:b3:scattering-theory:11}
(a) $a = \qty{5.3}{nm}$：$\sigma = 8\pi a^2 = \qty{7e-16}{m^2}$。
(b) $k = \sqrt{2mk_{\text{B}}T}/\hbar \approx \qty{6e6}{m^{-1}}$：
$ka \approx 0.03 \ll 1$。(c) 气体可由理想
（$a \approx 0$）经强排斥、吸引，直至
不稳定——相互作用成为一个实验旋钮。(d) 扫过
共振时，原子两两配对形成分子；当 $a \to
\infty$ 时，气体的强耦合程度堪比中子星物质，
只是以桌面形式呈现：借助\omterm{thm:b3:scattering-theory:partial}{散射长度}做量子
模拟。
\end{solution}

\begin{solution}{exo:b3:scattering-theory:12}
(a) $F(0) = \int\rho/Q = 1$。(b) 展开指数，
线性项平均为零，二次项给出
$-q^2\langle r^2\rangle/6$。(c) μ子氢的兰姆位移
（\cref{pb:b3:hydrogen-atom:1}）用完全不同的方法
测定同一个 $\langle r^2\rangle$；两者
在 $4\%$ 的水平上不一致（0.877 对 \qty{0.841}{fm}），
意味着要么存在微妙的系统误差，要么有新物理——于是
又花了十年重测（最终结果表明，主要是系统误差）。
(d) 质子是一团柔软的云，内部含有坚硬的颗粒：夸克
——卢瑟福的论证，再深一层。
\end{solution}

\begin{solution}{pb:b3:scattering-theory:1}
\textbf{1.} 每个原子的 $\theta \sim \num{8e-4}$ rad；即便是
穿过 $10^4$ 个原子层的无规行走，也只累积几度的偏转：
葡萄干布丁模型中后向散射不可能发生。
\textbf{2.} $E = k_CzZe^2/d_0$：$d_0 = k_CzZe^2/E =
\qty{4.1e-14}{m}$——即四十飞米。
\textbf{3.} 全部正电荷位于
$\sim\qty{4e-14}{m}$ 之内：比原子小一万倍。
\textbf{4.} $b \to 0$：$\theta \to 180^\circ$；$b \to \infty$：
$\theta \to 0$——正面碰撞反弹，远距离擦过。
\textbf{5.} $\theta > 90^\circ$ 需要 $b < d_0/2 =
\qty{2.1e-14}{m}$：比例为 $nt\,\pi b^2 \approx \num{3e-5}$——
对这块箔，平均每三万个 α 粒子中有一个。
\textbf{6.} 与盖革和马斯登的结果同量级（他们的
箔更厚，每八千个中有一个）：“不可能”的反弹
出现的频率，恰好与点状原子核所要求的一致。
\textbf{7.} $\dd\sigma = 2\pi b|\dd b|$，其中 $b =
(d_0/2)\cot(\theta/2)$：$|\dd b/\dd\theta| = (d_0/4)/\sin^2(
\theta/2)$，且 $\dd\Omega = 2\pi\sin\theta\,\dd\theta$；
合并时，利用 $\sin\theta = 2\sin(\theta/2)\cos(\theta/2)$
进行约化，得 $(d_0/4)^2\sin^{-4}(\theta/2)$。
\textbf{8.} $d_0/4 = k_CzZe^2/4E = Q_1Q_2/16\pi\varepsilon_0E$：
即玻恩结果——$1/r$ 的巧合。
\textbf{9.} $(d_0/4)^2 = \qty{1.06}{b/sr}$：$236$、$17$ 和
$\qty{1.9}{b/sr}$，分别对应 $30^\circ$、$60^\circ$、$120^\circ$。
\textbf{10.} $(\sin60^\circ/\sin15^\circ)^4 \approx 125$：他们的
计数跟随这类比值，跨越五个数量级的速率——这是
定律，而非趋势。
\textbf{11.} 在几何条件固定时，计数率正比于 $Z^2$：
比较不同的箔可以标定原子核电荷本身，为
将 $Z$ 等同于原子序数提供依据。
\textbf{12.} 所有计数率都降为四分之一；角度的\emph{形状}
保持不变——这是 $1/E^2$ 的清晰实验标志。
\textbf{13.} $d(180^\circ) = d_0 = \qty{41}{fm}$；$d(60^\circ) =
(d_0/2)(1 + 2) = \qty{62}{fm}$。
\textbf{14.} 该公式在最大角度处仍然成立：金
原子核小于 $\sim\qty{4e-14}{m}$。
\textbf{15.} 入射粒子开始触及原子核：
短程的\emph{强}相互作用（以及吸收）在
飞米量级的距离上开始起作用——偏离量度量了
原子核的边界。
\textbf{16.} $R = 1.2 \times 197^{1/3} = \qty{7.0}{fm}$——
明显小于卢瑟福给出的上界。
\textbf{17.} 体积 $\propto A$：核物质具有固定的
密度（$\sim\qty{2e17}{kg/m^3}$）——原子核是
不可压缩液体的液滴，这正是
\cref{ch:b3:nuclear-physics}的出发图像。
\textbf{18.} 光谱学探测的是电子；只有
能穿透原子的入射粒子，才能揭示
\emph{正}电荷与质量位于何处。
\textbf{19.} $20$ GeV 的电子以难以置信的
大角度发生散射并损失能量——对于一个
柔软的质子而言硬事件太多：卢瑟福式反弹的重演。
\textbf{20.} 电子不感受强相互作用，也没有已知的
内部结构：它是点状探针，只读取电荷，没有
α 粒子自身的复合性。
\textbf{21.} 弹性形状因子的骤降表明质子
是一团弥散的云；而硬的\emph{非弹性}
散射的持续出现说明云中含有点状成分——
夸克（部分子）。
\textbf{22.} 1911 年：$\sim\qty{e-14}{m}$；SLAC：$\hbar c/q \sim
\qty{0.2}{fm} = \qty{2e-16}{m}$；LHC：$\sim\qty{2e-19}{m}$——
一百年间，“标尺”缩小了五个数量级。
\textbf{23.} 氢原子的位置（在冰、聚合物、蛋白质中）以及
磁结构（反铁磁体、自旋螺旋）——两者
对 X 射线几乎不可见，而都是中子的拿手好戏。
\textbf{24.} 值得发现的过程，每
$10^{15}$ 次普通碰撞才出现一次：只有巨大的
亮度才能把飞靶换算为每年的事例数——对撞机设计
就是\omterm{def:b3:scattering-theory:cross-section}{截面}的算术。
\textbf{25.} 一条 $\sin^{-4}(\theta/2)$ 的定律，逐个闪烁核对，
将原子的质量集中于其体积的 $10^{-14}$ 之内；在
越来越高的 $q$ 下运行，同一实验测出了原子核的大小，
继而测出了其组成成分，至今仍在探索。
\end{solution}
